\section{Introduction}\label{sec:introduction}

How much work does it take to optimize a linear function subject to linear
inequalities? For rational coefficients, the ellipsoid method of
Khachiyan~\cite{Khach,Khach2} gave an answer polynomial in the length of the
input. Two linear programs with the same numbers of variables and constraints
can have very different input lengths. Megiddo~\cite{MegEnc,Meg} asked
whether those two dimensions alone could bound the number of arithmetic
operations. Tarjan raised the question for real data~\cite[pp.~1620--1621]{Tari},
and Smale made it his ninth problem~\cite{Smale}. We give an affirmative answer.

We maximize \(c^\top x\) subject to \(Ax\le b\), given
\(A\in\RR^{m\times n}\), \(b\in\RR^m\) and \(c\in\RR^n\) explicitly.
A \emph{strongly polynomial} algorithm takes polynomially many steps
in \(m,n\), counting exact scalar arithmetic, comparisons and control work.
On rational input every intermediate number has bit length polynomial in
the full binary input length.

\begin{theorem}[Strongly polynomial linear programming]\label{thm:lp}
For arbitrary real \(A,b,c\), a deterministic strongly polynomial algorithm
returns one of the following certificates for \(\max\{c^\top x:Ax\le b\}\).
\begin{enumerate}
 \item \emph{Finite optimum.} A strictly complementary primal--dual pair \(x,y\) with
 \[
    s=b-Ax\ge0,\qquad y\ge0,\qquad A^\top y=c,\qquad
    s^\top y=0,\qquad s+y>0 .
 \]
 \item \emph{Unbounded objective.} A feasible point \(x\) and an improving ray
 \(r\) with \(Ar\le0\) and \(c^\top r>0\).
 \item \emph{Infeasibility.} A vector \(y\ge0\) with \(A^\top y=0\),
 \(b^\top y<0\), and \(\lvert\supp y\rvert\le 1+\rank A\).
\end{enumerate}
\end{theorem}

We use a uniform sequential model~\cite{BSS}. The scalar operations are
addition, subtraction, multiplication, division by a nonzero scalar,
comparison and copying, with constants \(0,1\). Control operations on
indices and explicit bit strings are charged at bit cost. The algorithm accesses scalar inputs only through the scalar operations.
All vector inequalities are componentwise. On rational data we measure a
number by its reduced fraction.

The arithmetic result extends to every ordered field. In particular, on
rational input the certificates are exact rationals. If the input is supplied
as integer numerator--denominator pairs, an implementation also keeps the
integers in its unreduced pairs of polynomial bit length throughout the
computation (Section~\ref{sec:arithmetic}).

\paragraph{How the algorithm finds a certificate.}
Consider first the feasibility of \(Ax\le b\). A feasible point has
nonnegative slacks \(b-Ax\). A certificate of infeasibility instead gives
nonnegative weights for a combination of the inequalities that yields a
contradiction. Homogenizing the slack vectors puts these alternatives in
two orthogonal subspaces. To decide feasibility it suffices to determine
which subspace contains a nonnegative vector positive at one designated
coordinate.

A positive entry of such a vector can be arbitrarily small compared with its
largest entry. Changing the unit of a coordinate changes which entries can be
seen at a fixed precision. The algorithm uses these changes of unit as
queries. At each query it obtains reliable signs where the rescaled entries
are largest. A search combining the signs from different queries identifies
the subspace containing a witness, even though the algorithm does not know
in advance which signs it can trust. It then sets coordinates to zero, one
at a time whenever a witness survives, until linear algebra recovers the
witness exactly.

The difficulty is making each query affordable. For a witness of minimal support,
the ratios of its nonzero entries are ratios of subdeterminant magnitudes
of a basis matrix for the two subspaces. Even in a fixed dimension, the
logarithms of these magnitudes can be arbitrarily far apart.
The algorithm gives short substitutes for the logarithms as they are
needed, preserving the comparisons that determine which entries are large.
This lets it compute each set of signs using ordinary polynomial-time LP
on rational data whose bit lengths are bounded in the dimensions alone.
Section~\ref{sec:overview} explains how the search, the changes of unit and
the compression fit together.

\subsection{Consequences}

The same search applies to a convex quadratic objective under linear
constraints. For arbitrary real coefficients we find an exact minimizer
and its multipliers, or certify infeasibility or unboundedness, in
deterministic strongly polynomial time (Corollary~\ref{cor:qp}).
The underlying result solves the monotone linear complementarity problem.
Given \(M,q\) with \(M+M^\top\) positive semidefinite, it finds \(z\ge0\)
such that \(Mz+q\ge0\) and \(z^\top(Mz+q)=0\), or certifies that no
such \(z\) exists (Theorem~\ref{thm:lcp}).

There is a second use for the signs returned by a query. If a point differs
from the fixed point of a contraction in max norm, the map moves it toward
that fixed point in every coordinate of greatest error. For a self-map of
\([-1,1]^n\) with Lipschitz constant at most \(1-2^{-b}\), this gives a
deterministic algorithm that locates the fixed point to within \(2^{-k}\)
in time polynomial in \(n,b,k\), counting all bits exchanged with the
evaluation oracle (Theorem~\ref{thm:contraction}).

In a finite turn-based mean-payoff game, two players with perfect information
choose moves on a graph. One maximizes and the other minimizes the long-run
average reward. The fixed-point algorithm and policy evaluation by LP give the exact values
from every starting vertex and optimal strategies for both players. The running time
is strongly polynomial when each move has a deterministic outcome and
polynomial in the binary input length when outcomes are stochastic and the
data rational. This includes polynomial time for simple stochastic reachability
games. Theorem~\ref{thm:games} gives the bounds for arbitrary real rewards
and for discounted play.

\subsection{Earlier approaches to strongly polynomial time}
\label{sec:lp-context}

The first polynomial-time algorithms for general linear programming
measure the size of a rational instance in bits. Khachiyan's ellipsoid
algorithm~\cite{Khach,Khach2} draws on the cutting-plane methods of
Shor and Yudin--Nemirovskii~\cite{Shor,YN}. In its feasibility argument,
the volume of the enclosing ellipsoid falls by a controlled factor at
each cut. Determinant bounds in terms of the input length tell the
algorithm when that volume is too small to contain a feasible region
of the required thickness~\cite{GLS}. Karmarkar's projective
algorithm~\cite{Kar} gave a second route to polynomial bit complexity.
Subsequent interior-point methods follow solutions of penalized problems
as the penalty for approaching the boundary decreases~\cite{NN}.
Faster linear algebra reduces the work per solve~\cite{Ren,Vai,CLS,VDB},
while exact recovery still has to account for the input
precision~\cite{GPV}. Even on the interval $[0,1]$, the objective
$x\mapsto 2^{-L}x$ gives the two vertices values only $2^{-L}$ apart.
Intervals of radius $2^{-k}$ around these values are disjoint only for
$k>L+1$. A single exact sign comparison selects the optimal endpoint.

\paragraph{Removing dependence on the data.}
Tardos first obtained a strongly polynomial algorithm for minimum-cost
circulation~\cite{TarCirc}, then showed that the arithmetic operation
count for a rational LP can depend on the encoding length of the
constraint matrix alone~\cite{Tardos}. This gives a strongly polynomial
bound whenever the entries of the matrix have polynomial encoding length
in the dimensions, including matrices with entries in $\{0,1,-1\}$.
Frank and Tardos~\cite{FT} replace a vector by a short integer vector
whose inner products with bounded integer vectors have the same signs.
When the constraint matrix has short entries, this lets one compress
the objective and the right-hand side while preserving the comparisons
needed for optimization~\cite[\S6.6]{GLS}.

A different source of strongly polynomial bounds is structure in the
constraint matrix. Megiddo solved LP in linear time for fixed
dimension~\cite{MegLin}, and solved feasibility for inequalities
involving at most two variables each~\cite{Meg}. Generalized flow
maximization supplies a related class with at most two nonzeros per
column~\cite{Vegh,OV}. Dadush, Koh, Natura, Olver and V\'egh obtained a
strongly polynomial optimization algorithm for every matrix with at
most two nonzeros in each row or at most two in each
column~\cite{DKNOV}. Their later work treats classes for which the
central path has polynomial straight-line complexity~\cite{DKNOV26}.

\paragraph{Geometry of the constraint matrix.}
For a real matrix $A$, Vavasis and Ye~\cite{VY} gave a path-following
algorithm whose iteration bound depends on the dimensions and the
logarithm of a condition measure of $A$. Their layered least-squares
steps can cross several numerical scales at once. A related condition
measure in standard form $Ax=b$, $x\ge0$ has a concrete description.
A \emph{circuit} of $A$ is a nonzero kernel vector of minimal support.
The \emph{circuit imbalance} is the largest ratio of magnitudes of
two nonzero coordinates in the same circuit~\cite{Rock,CIRC}.
Dadush, Huiberts, Natura and V\'egh~\cite{DHNV} obtained an iteration
bound in terms of the best condition achievable by diagonal scaling of
the columns. Circuit imbalance also controls the precision needed in
an approximate-to-exact solver~\cite{DNV}.

Even the best rescaling can leave large ratios in a fixed dimension.
For $0<u<1$, the matrix
\[
 A=\begin{pmatrix}1&0&1&u\\0&1&u&1\end{pmatrix}
\]
has circuits $(-1,-u,1,0)$ and $(-u,-1,0,1)$. The ratios of the
magnitudes of their first two coordinates are $1/u$ and $u$.
Any positive diagonal column scaling multiplies these two ratios by the
same factor, so at least one ratio or its reciprocal is at least $1/u$.
A separate obstruction comes from the central path itself. Allamigeon,
Benchimol, Gaubert and Joswig constructed instances requiring
exponentially many segments to follow the logarithmic central path
in a prescribed neighborhood~\cite{ABGJ}. Allamigeon, Gaubert and
Vandame extended this obstruction to path-following with
self-concordant barriers~\cite{AGV}.

\paragraph{Edge paths and rescalings.}
Early applications ranged from Kantorovich's production planning~\cite{Kan}
to the computation of optimal strategies in zero-sum matrix games by LP
duality~\cite{Gale}. Dantzig's simplex method exploits the geometry of these
problems by moving between adjacent feasible bases~\cite{Dantzig}.
Klee and Minty exhibited exponentially long runs for a classical
choice of pivot~\cite{KM}. Randomized pivot rules achieve
subexponential bounds~\cite{Kalai,MSW}, and smoothed analysis gives
polynomial expected bounds in the dimensions and the inverse noise
scale~\cite{ST,BH}. Awoniyi has proposed a different pivot-based
approach to strongly polynomial LP~\cite{Awon}.

Even the shortest possible edge walk can be long. Kalai and
Kleitman proved a quasipolynomial diameter bound for polyhedral
graphs~\cite{KK}, and Santos disproved the original Hirsch bound of
the number of facets minus the dimension~\cite{Santos}. The companion
paper~\cite{OAIHirsch} constructs simple rational polytopes whose
graph diameter grows faster than any polynomial in their number of
facets. A circuit walk can leave the edge graph by taking a maximal
feasible step in a circuit direction. Natura constructs polynomial-length
monotone circuit walks between specified vertices~\cite{Natura}.

The algorithm of Theorem~\ref{thm:lp} searches over coordinate
rescalings. In earlier rescaling algorithms, finding that one coordinate
must be small in every normalized feasible point can guide a change
of scale that improves the geometry~\cite{Betke,DV,Chu,DVZ,PS}. Here a
rescaling selects the coordinates on which a primal or dual witness
can be detected. Section~\ref{sec:overview} explains how the algorithm
combines such observations, even when the initial system has arbitrarily
large circuit imbalance.

